\subsection{Fórmula de las dimensiones y anillos universalmente catenarios.}
\label{IV.2.s5.6-fr}

\begin{proposition}[5.6.1]
\EGAReaderTarget{ega4.c04.prop.5-6-1}{IV.2.s5.6.1-fr}
Sea $A$ un anillo noetheriano. Las propiedades siguientes son equivalentes:
\begin{enumerate}
\item[\textup{a)}] Todo anillo de polinomios $A[T_1,\ldots,T_n]$ es catenario.
\item[\textup{b)}] Toda álgebra de tipo finito sobre $A$ es catenaria.
\item[\textup{c)}] $A$ es catenario, y para toda $A$-álgebra íntegra local,
esencialmente de tipo finito \textup{(1.3.8)}, $B'$, se tiene, denotando por
$\mathfrak p$ la imagen inversa en $A$ del ideal maximal $\mathfrak q$ de
$B'$, por $A'$ la imagen de $A_{\mathfrak p}$ en $B'$, y por $k$ y $k'$ los
cuerpos residuales de $A'$ y $B'$, respectivamente,
\begin{equation}
\dim(A')+\operatorname{deg.tr}_{A'}(B')
=\dim(B')+\operatorname{deg.tr}_k k'.
\tag{5.6.1.1}\label{IV.2.s5.6.1.1-fr}
\end{equation}
\end{enumerate}
\end{proposition}

El hecho de que \emph{b)} implica \emph{c)} se sigue de (5.5.8); en efecto, $B'$ es una
$A'$-álgebra local esencialmente de tipo finito (1.3.10), por tanto de la
forma $B''_{\mathfrak q''}$, donde $B''$ es una sub-$A'$-álgebra de tipo
finito de $B'$ y $\mathfrak q''$ un ideal primo de $B''$ situado sobre el
ideal maximal $\mathfrak p'$ de $A'$; además, los cuerpos de fracciones de
$B'$ y $B''$ son los mismos, por lo que
$\operatorname{deg.tr}_{A'}(B')=\operatorname{deg.tr}_{A'}(B'')$.
Para demostrar (5.6.1.1), basta mostrar que (bajo la hipótesis \emph{b)}) toda
sub-$A'$-álgebra $A_1$ de tipo finito de $B'[T]$ es catenaria; en efecto, los
dos miembros de (5.5.8.1), donde se sustituyen $A$, $B$ y $\mathfrak q$ por
$A'$, $B''$ y $\mathfrak q''$, serán entonces iguales, de donde se obtiene la
igualdad (5.6.1.1). Ahora bien, la hipótesis \emph{b)} implica que toda $A$-álgebra
esencialmente de tipo finito es catenaria (\textbf 0, 16.1.4), y $A_1$ es una
$A$-álgebra de este tipo (1.3.9).

Es trivial que \emph{b)} implica \emph{a)}; recíprocamente, \emph{a)} implica \emph{b)}, pues toda
$A$-álgebra de tipo finito es cociente de un álgebra de polinomios
(\textbf 0, 16.1.4).

Queda probar que \emph{c)} implica \emph{b)}. Como todo cociente por un ideal primo
\oldpage[IV-2]{98}
de una $A$-álgebra de tipo finito es una $A$-álgebra íntegra de tipo finito, se trata
de ver que, si $B$ es una $A$-álgebra íntegra de tipo finito y $\mathfrak q$ y
$\mathfrak q'$ son dos ideales primos de $B$ tales que
$\mathfrak q'\subset\mathfrak q$, se tiene (\textbf 0, 16.1.4.2)
\begin{equation}
\dim(B_{\mathfrak q}/\mathfrak q'B_{\mathfrak q})
+\dim(B_{\mathfrak q'})=\dim(B_{\mathfrak q}).
\tag{5.6.1.2}\label{IV.2.s5.6.1.2-fr}
\end{equation}

Sean $\mathfrak p$ y $\mathfrak p'$ las imágenes inversas de $\mathfrak q$ y
$\mathfrak q'$, respectivamente, en $A$, y $\mathfrak n$ el núcleo del
homomorfismo $A_{\mathfrak p}\to B_{\mathfrak q}$.

Como la imagen $A'$ de $A_{\mathfrak p}$ en $B_{\mathfrak q}$ es isomorfa a
$A_{\mathfrak p}/\mathfrak n$, la fórmula (5.6.1.1), aplicada a $A'$ y a
$B'=B_{\mathfrak q}$, se escribe
\begin{equation}
\dim(A_{\mathfrak p}/\mathfrak n)
+\operatorname{deg.tr}_{A_{\mathfrak p}/\mathfrak n}(B_{\mathfrak q})
=\dim(B_{\mathfrak q})
+\operatorname{deg.tr}_{\mathbf{k}(\mathfrak p)}\boldsymbol{k}(\mathfrak q).
\tag{5.6.1.3}\label{IV.2.s5.6.1.3-fr}
\end{equation}

Por otra parte, el núcleo de $A_{\mathfrak p'}\to B_{\mathfrak q'}$ es
$\mathfrak nA_{\mathfrak p'}$, por lo que, aplicando la fórmula (5.6.1.1),
donde se sustituye $A'$ por
$A_{\mathfrak p'}/\mathfrak nA_{\mathfrak p'}$ y $B'$ por
$B_{\mathfrak q'}$, se tiene
\begin{equation}
\dim(A_{\mathfrak p'}/\mathfrak nA_{\mathfrak p'})
+\operatorname{deg.tr}_{A_{\mathfrak p'}/\mathfrak nA_{\mathfrak p'}}
(B_{\mathfrak q'})
=\dim(B_{\mathfrak q'})
+\operatorname{deg.tr}_{\mathbf{k}(\mathfrak p')}\boldsymbol{k}(\mathfrak q').
\tag{5.6.1.4}\label{IV.2.s5.6.1.4-fr}
\end{equation}

Finalmente, $B/\mathfrak q'$ es una $A$-álgebra íntegra de tipo finito, la
imagen inversa de $\mathfrak q/\mathfrak q'$ en $A$ es $\mathfrak p$, y el
núcleo del homomorfismo
$A_{\mathfrak p}\to B_{\mathfrak q}/\mathfrak q'B_{\mathfrak q}$ es
$\mathfrak p'A_{\mathfrak p}$, por lo que, aplicando (5.6.1.1), donde se
sustituye $A'$ por $A_{\mathfrak p}/\mathfrak p'A_{\mathfrak p}$ y $B'$ por
$B_{\mathfrak q}/\mathfrak q'B_{\mathfrak q}$, se tiene
\begin{equation}
\dim(A_{\mathfrak p}/\mathfrak p'A_{\mathfrak p})
+\operatorname{deg.tr}_{A_{\mathfrak p}/\mathfrak p'A_{\mathfrak p}}
(B_{\mathfrak q}/\mathfrak q'B_{\mathfrak q})
=\dim(B_{\mathfrak q}/\mathfrak q'B_{\mathfrak q})
+\operatorname{deg.tr}_{\mathbf{k}(\mathfrak p)}\boldsymbol{k}(\mathfrak q).
\tag{5.6.1.5}\label{IV.2.s5.6.1.5-fr}
\end{equation}

Sumemos ahora (5.6.1.4) y (5.6.1.5) miembro a miembro, y observemos que
$\boldsymbol{k}(\mathfrak p')$ (resp. $\boldsymbol{k}(\mathfrak q')$) es el cuerpo de fracciones de
$A_{\mathfrak p}/\mathfrak p'A_{\mathfrak p}$ (resp.
$B_{\mathfrak q}/\mathfrak q'B_{\mathfrak q}$); por otra parte,
$A_{\mathfrak p'}/\mathfrak nA_{\mathfrak p'}$ y
$A_{\mathfrak p}/\mathfrak n$ tienen el mismo cuerpo de fracciones, y
$B_{\mathfrak q'}$ y $B_{\mathfrak q}$ tienen el mismo cuerpo de fracciones;
finalmente, como se supone que $A$ es catenario, se tiene (\textbf 0, 16.1.4.2)
\[
\dim(A_{\mathfrak p}/\mathfrak p'A_{\mathfrak p})
+\dim(A_{\mathfrak p'}/\mathfrak nA_{\mathfrak p'})
=\dim(A_{\mathfrak p}/\mathfrak n)
\]
Teniendo en cuenta estas observaciones y (5.6.1.3), se obtiene la relación
(5.6.1.2). C.Q.F.D.

